\begin{figure}[htbp]
			\centering
			\resizebox{\textwidth}{!}{%
				\begin{tikzpicture}[>=stealth, font=\small,
					asta/.style    = {line width=1.2pt, line cap=round, line join=round, black},
					nodo/.style    = {circle, fill=white, draw=black,
						inner sep=0pt, minimum size=6pt, line width=1.2pt},
					etasta/.style  = {circle, draw=black, fill=white,
						inner sep=1pt, font=\scriptsize, line width=0.6pt},
					quota/.style   = {<->, >={Latex[length=1.6mm]}, line width=0.5pt},
					ext/.style     = {line width=0.3pt},
					lab/.style     = {font=\small},
					]
					
					%% ── Parametri geometrici globali ─────────────────────────────────────────
					\def\L{2.0}       % lunghezza ell (asta di corrente)
					\def\h{1.6}       % altezza tra correnti
					\def\DX{9.5}      % offset orizzontale tra pannelli
					\def\DY{-7.0}     % offset verticale tra pannelli
					
					%% Posizioni nodi (origine = N1 di ciascun pannello)
					\pgfmathsetmacro{\xNa}{0}           % N1
					\pgfmathsetmacro{\xNb}{\L}          % N2
					\pgfmathsetmacro{\xNc}{2*\L}        % N3
					\pgfmathsetmacro{\xNd}{3*\L}        % N4
					\pgfmathsetmacro{\xNe}{0.5*\L}      % N5
					\pgfmathsetmacro{\xNf}{1.5*\L}      % N6
					\pgfmathsetmacro{\xNg}{2.5*\L}      % N7
					
					%% =========================================================================
					%% PANNELLO 1 — STRUTTURA QUOTATA + CARICO UNITARIO MOBILE
					%% =========================================================================
					\begin{scope}[xshift=0cm, yshift=0cm]
						
						\coordinate (N1) at (\xNa, 0);
						\coordinate (N2) at (\xNb, 0);
						\coordinate (N3) at (\xNc, 0);
						\coordinate (N4) at (\xNd, 0);
						\coordinate (N5) at (\xNe, \h);
						\coordinate (N6) at (\xNf, \h);
						\coordinate (N7) at (\xNg, \h);
						
						%% --- aste
						\draw[asta](N1)--(N5);
						\draw[asta](N2)--(N5);
						\draw[asta](N2)--(N6);
						\draw[asta](N3)--(N6);
						\draw[asta](N3)--(N7);
						\draw[asta](N4)--(N7);
						\draw[asta](N1)--(N2)--(N3)--(N4);
						\draw[asta](N5)--(N6)--(N7);
						
						%% --- cerniera in N1
						\draw[line width=1.0pt](N1)--++(-0.24,-0.40)--++(0.48,0)--cycle;
						\draw[line width=1.0pt]($(N1)+(-0.34,-0.40)$)--($(N1)+(0.34,-0.40)$);
						\foreach \x in {-0.28,-0.14,0.00,0.14,0.28}
						\draw[very thin]($(N1)+(\x,-0.40)$)--++(-.10,-0.13);
						
						%% --- carrello in N4
						\draw[line width=1.0pt](N4)--++(-0.24,-0.34)--++(0.48,0)--cycle;
						\draw[line width=1.0pt]($(N4)+(-0.34,-0.34)$)--($(N4)+(0.34,-0.34)$);
						\draw[line width=1.0pt]($(N4)+(-0.16,-0.47)$) circle(2.2pt);
						\draw[line width=1.0pt]($(N4)+( 0.16,-0.47)$) circle(2.2pt);
						\draw[line width=1.0pt]($(N4)+(-0.34,-0.58)$)--($(N4)+(0.34,-0.58)$);
						\foreach \x in {-0.28,-0.14,0.00,0.14,0.28}
						\draw[very thin]($(N4)+(\x,-0.58)$)--++(-.10,-0.13);
						
						%% --- pallini dei nodi
						\foreach \n in {N1,N2,N3,N4,N5,N6,N7}
						{\node[nodo]at(\n){};}
						
						%% --- etichette dei nodi
						\node[left,       lab] at (N1){$1$};
						\node[below,      lab] at (N2){$2$};
						\node[below,      lab] at (N3){$3$};
						\node[right,      lab] at (N4){$4$};
						\node[above left, lab] at (N5){$5$};
						\node[above,      lab] at (N6){$6$};
						\node[above right,lab] at (N7){$7$};
						
						%% --- etichette delle aste (cerchiettate)
						\node[etasta] at (0.5*\L, -0.28)   {$1$};
						\node[etasta] at (1.5*\L, -0.28)   {$2$};
						\node[etasta] at (2.5*\L, -0.28)   {$3$};
						\node[etasta] at (\L,     \h+0.28) {$4$};
						\node[etasta] at (2*\L,   \h+0.28) {$5$};
						\node[etasta] at (0.15*\L, 0.5*\h+0.10){$6$};
						\node[etasta] at (0.60*\L, 0.5*\h+0.10){$7$};
						\node[etasta] at (1.15*\L, 0.5*\h+0.10){$8$};
						\node[etasta] at (1.60*\L, 0.5*\h+0.10){$9$};
						\node[etasta] at (2.15*\L, 0.5*\h+0.10){$10$};
						\node[etasta] at (2.85*\L, 0.5*\h+0.10){$11$};
						
						%% --- carico unitario mobile sul CORRENTE INFERIORE
						\pgfmathsetmacro{\xPi}{1.4*\L}       % punto di applicazione inferiore
						\def\Fh{0.70}                         % altezza della freccia
						\draw[->, >={Latex[length=1.8mm]}, line width=0.8pt]
						(\xPi, \Fh) -- (\xPi, 0.05);
						\node[right, font=\footnotesize] at (\xPi-0.1, \Fh-0.3){$F=1$};
						%\node[nodo, minimum size=3pt] at (\xPi, 0){};
						
						%% --- carico unitario mobile sul CORRENTE SUPERIORE
						\pgfmathsetmacro{\xPs}{0.75*\L}            % punto di applicazione superiore
						\draw[->, >={Latex[length=1.8mm]}, line width=0.8pt]
						(\xPs, \h+\Fh) -- (\xPs, \h+0.05);
						\node[above, font=\footnotesize] at (\xPs, \h+\Fh-0.1){$F=1$};
						%\node[nodo, minimum size=3pt] at (\xPs, \h){};
						
						%% --- quote orizzontali: tre segmenti di ell
						\def\yq{-1.25}
						\draw[ext](\xNa, -0.78)--(\xNa, \yq-0.10);
						\draw[ext](\xNb,  0)   --(\xNb, \yq-0.10);
						\draw[ext](\xNc,  0)   --(\xNc, \yq-0.10);
						\draw[ext](\xNd, -0.78)--(\xNd, \yq-0.10);
						\draw[quota](\xNa,\yq)--(\xNb,\yq) node[midway,below,lab]{$\ell$};
						\draw[quota](\xNb,\yq)--(\xNc,\yq) node[midway,below,lab]{$\ell$};
						\draw[quota](\xNc,\yq)--(\xNd,\yq) node[midway,below,lab]{$\ell$};
						
						%% --- quota ascissa z del carico inferiore (sotto, su linea intermedia)
						\def\yz{-0.85}
						\draw[ext](\xNa, 0)  --(\xNa, \yz-0.10);
						\draw[ext](\xPi, 0)  --(\xPi, \yz-0.10);
						\draw[quota](\xNa, \yz)--(\xPi, \yz) node[midway,above,font=\footnotesize]{$z$};
						
						%% --- quota ascissa z' del carico superiore (sopra la travatura)
						\def\yzp{\h+1.25}
						\draw[ext](\xNa, 0)    --(\xNa, \yzp+0.10);
						\draw[ext](\xPs, \h+\Fh+0.25)--(\xPs, \yzp+0.10);
						\draw[quota](\xNa, \yzp)--(\xPs, \yzp) node[midway,below,font=\footnotesize]{$z'$};
						
						%% --- quota verticale h
						\def\xq{-1.00}
						\draw[ext](-0.48, 0)--(\xq-0.10, 0);
						\draw[ext](N5)      --(\xq-0.10, \h);
						\draw[quota](\xq,0)--(\xq,\h) node[midway,left,lab]{$h$};
						
						\node[font=\small, align=center] at (1.5*\L, -2.50)
						{(a) Struttura con carico mobile unitario};
						
					\end{scope}
					
					%% =========================================================================
					%% PANNELLO 2 — MECCANISMO VIRTUALE: DUE TRONCONI, DISTORSIONE SU ASTA 4,
					%%              CAMPI v* (sotto la travatura) E u* (a destra della travatura)
					%% =========================================================================
					\begin{scope}[xshift=\DX cm, yshift=0cm]
						
						%% --- coordinate dei nodi
						\coordinate (N1) at (\xNa, 0);
						\coordinate (N2) at (\xNb, 0);
						\coordinate (N3) at (\xNc, 0);
						\coordinate (N4) at (\xNd, 0);
						\coordinate (N5) at (\xNe, \h);
						\coordinate (N6) at (\xNf, \h);
						\coordinate (N7) at (\xNg, \h);
						
						%% --- riempimenti grigi dei due tronconi
						\fill[black!12] (N1) -- (N2) -- (N5) -- cycle;
						\fill[black!12] (N2) -- (N3) -- (N4) -- (N7) -- (N6) -- cycle;
						
						%% --- aste (asta 4, N5--N6, in grigio chiaro)
						\draw[asta](N1)--(N5);
						\draw[asta](N2)--(N5);
						\draw[asta](N2)--(N6);
						\draw[asta](N3)--(N6);
						\draw[asta](N3)--(N7);
						\draw[asta](N4)--(N7);
						\draw[asta](N1)--(N2)--(N3)--(N4);
						\draw[asta, gray!55](N5)--(N6);      % asta 4: cede (grigia)
						\draw[asta](N6)--(N7);
						
						
						%% --- distorsione sull'asta 4: due frecce antagoniste sopra l'asta
						\pgfmathsetmacro{\yf}{\h+0.30}
						\pgfmathsetmacro{\xfc}{0.5*(\xNe+\xNf)}
						\def\frL{0.30}                          % lunghezza di ciascuna freccia
						\draw[->, >={Latex[length=1.6mm]}, line width=0.8pt]
						(\xfc-0.06, \yf) -- (\xfc-\frL, \yf);
						\draw[->, >={Latex[length=1.6mm]}, line width=0.8pt]
						(\xfc+0.06, \yf) -- (\xfc+\frL, \yf);
						\node[above, font=\footnotesize] at (\xfc, \yf+0.0){$\varrho_w=1$};
						
						%% --- cerniera in N1
						\draw[line width=1.0pt](N1)--++(-0.24,-0.40)--++(0.48,0)--cycle;
						\draw[line width=1.0pt]($(N1)+(-0.34,-0.40)$)--($(N1)+(0.34,-0.40)$);
						\foreach \x in {-0.28,-0.14,0.00,0.14,0.28}
						\draw[very thin]($(N1)+(\x,-0.40)$)--++(-.10,-0.13);
						
						%% --- carrello in N4
						\draw[line width=1.0pt](N4)--++(-0.24,-0.34)--++(0.48,0)--cycle;
						\draw[line width=1.0pt]($(N4)+(-0.34,-0.34)$)--($(N4)+(0.34,-0.34)$);
						\draw[line width=1.0pt]($(N4)+(-0.16,-0.47)$) circle(2.2pt);
						\draw[line width=1.0pt]($(N4)+( 0.16,-0.47)$) circle(2.2pt);
						\draw[line width=1.0pt]($(N4)+(-0.34,-0.58)$)--($(N4)+(0.34,-0.58)$);
						\foreach \x in {-0.28,-0.14,0.00,0.14,0.28}
						\draw[very thin]($(N4)+(\x,-0.58)$)--++(-.10,-0.13);
						
						%% --- pallini dei nodi
						\foreach \n in {N1,N2,N3,N4,N5,N6,N7}
						{\node[nodo]at(\n){};}
						
						%% --- etichette dei centri sulla struttura
						\node[above left,  font=\footnotesize] at (N1){$C_L$};
						\node[below,       font=\footnotesize] at (N2){$C_{L,R}$};
						\node[above right, font=\footnotesize] at (N4){$C_R$};
						
						%% ═════════════════════════════════════════════════════════════════
						%% RETTA DI PROIEZIONE ORIZZONTALE — campo v* (sotto la travatura)
						%% ═════════════════════════════════════════════════════════════════
						\def\yB{-2.30}
						\def\vpk{0.80}
						
						\draw[ext, dashed, gray!70](\xNa, -0.78)--(\xNa, \yB);
						\draw[ext, dashed, gray!70](\xNb,  0   )--(\xNb, \yB);
						\draw[ext, dashed, gray!70](\xNd, -0.78)--(\xNd, \yB);
						
						\fill[black!12]
						(\xNa, \yB) -- (\xNb, \yB+\vpk) -- (\xNd, \yB) -- cycle;
						\draw[line width=1.0pt]
						(\xNa, \yB) -- (\xNb, \yB+\vpk) -- (\xNd, \yB);
						\draw[line width=0.4pt](\xNa-0.40, \yB)--(\xNd+0.40, \yB);
						
						\foreach \x in {\xNa, \xNb, \xNd}
						\draw[line width=0.5pt](\x, \yB-0.08)--(\x, \yB+0.08);
						
						%% --- archetti degli angoli ϑ_L e ϑ_R al piede della spezzata v*
						\def\Rav{0.55}
						
						%% ϑ_L in C_L
						\pgfmathsetmacro{\angLv}{atan2(\vpk, \xNb-\xNa)}
						\draw[line width=0.5pt]
						(\xNa, \yB) ++(0:\Rav) arc (0:\angLv:\Rav);
						\node[font=\footnotesize]
						at ($(\xNa, \yB) + ({0.5*\angLv}:\Rav) + (0.45, 0.05)$){$\vartheta_L^*$};
						%%                                       ^^^^^^^^^^^^^
						%%                                       offset (dx, dy) dell'etichetta
						
						%% ϑ_R in C_R
						\pgfmathsetmacro{\angRv}{180 - atan2(\vpk, \xNd-\xNb)}
						\draw[line width=0.5pt]
						(\xNd, \yB) ++(\angRv:\Rav) arc (\angRv:180:\Rav);
						\node[font=\footnotesize]
						at ($(\xNd, \yB) + ({0.5*(\angRv+180)}:\Rav) + (-0.25, -0.25)$){$\vartheta_R^*$};
						%%                                             ^^^^^^^^^^^^^
						
						\node[below, font=\footnotesize] at (\xNa, \yB-0.08){$C_L$};
						\node[below, font=\footnotesize] at (\xNb, \yB-0.08){$C_{L,R}$};
						\node[below, font=\footnotesize] at (\xNd, \yB-0.08){$C_R$};
						
						\node[above right=-1pt, font=\footnotesize]
						at (\xNb, \yB+\vpk){$\dfrac{2}{3}\,\dfrac{\varrho_w\,\ell}{h}$};
						%\node[left, font=\small] at (\xNa-0.45, \yB+\vpk/2){$v^*$};
						
						%% etichette dei due campi v*_L e v*_R sotto la baseline
						\node[below, font=\footnotesize] at ($(\xNa, \yB)!0.5!(\xNb, \yB) + (0, +0.98)$){$v^*_L$};
						\node[below, font=\footnotesize] at ($(\xNb, \yB)!0.5!(\xNd, \yB) + (0, +0.98)$){$v^*_R$};
						
						
						
						%% ═════════════════════════════════════════════════════════════════
						%% RETTA DI PROIEZIONE VERTICALE — campo u* (a destra della travatura)
						%% ═════════════════════════════════════════════════════════════════
						\pgfmathsetmacro{\xU}{\xNd+1.40}
						\def\dL{0.5333}
						\def\dR{0.2667}
						
						%% linee di richiamo orizzontali dalla struttura alla retta u*
						\draw[ext, dashed, gray!70](\xNd+0.40, 0) -- (\xU+0.10, 0);
						\draw[ext, dashed, gray!70](\xNg+0.10, \h) -- (\xU+0.10, \h);
						
						%% retta verticale non ruotata (riferimento tratteggiato)
						\draw[ext, dashed, gray!70](\xU, 0) -- (\xU, \h);
						
						%% piede: tick coi tre centri sovrapposti
						\draw[line width=0.5pt](\xU-0.08, 0) -- (\xU+0.08, 0);
						\node[below=1pt, font=\footnotesize] at (\xU, -0.05)
						{$C_L\!\equiv\!C_{L,R}\!\equiv\!C_R$};
						
						%% retta ruotata CCW (troncone sinistro)
						\draw[line width=0.9pt](\xU, 0) -- (\xU-\dL, \h);
						%% retta ruotata CW (troncone destro)
						\draw[line width=0.9pt](\xU, 0) -- (\xU+\dR, \h);
						
						%				%% --- archetti degli angoli ϑ_L e ϑ_R al piede della retta u*
						%				\def\Ra{0.55}
						%				\pgfmathsetmacro{\angL}{90 + atan2(\dL, \h)}
						%				\draw[line width=0.5pt]
						%				(\xU, 0) ++(90:\Ra) arc (90:\angL:\Ra);
						%				\node[above, font=\footnotesize]
						%				at ($(\xU, 0) + ({0.5*(90+\angL)}:\Ra+0.20)$){$\vartheta_L^*$};
						%				
						%				\pgfmathsetmacro{\angR}{90 - atan2(\dR, \h)}
						%				\draw[line width=0.5pt]
						%				(\xU, 0) ++(\angR:\Ra) arc (\angR:90:\Ra);
						%				\node[above, font=\footnotesize]
						%				at ($(\xU, 0) + ({0.5*(90+\angR)}:\Ra+0.20)$){$\vartheta_R^*$};
						
						%% quota del gap ϱ_w al livello h
						\draw[ext](\xU-\dL, \h) -- (\xU-\dL, \h+0.30);
						\draw[ext](\xU+\dR, \h) -- (\xU+\dR, \h+0.30);
						\draw[quota]
						(\xU-\dL, \h+0.25) -- (\xU+\dR, \h+0.25)
						node[midway, above, font=\footnotesize]{$\varrho_w$};
						
						%\node[right=2pt, font=\small] at (\xU+\dR, \h/2){$u^*$};
						%% etichette dei due campi u*_L e u*_R
						\node[font=\footnotesize] at ($(\xU-\dL, \h) + (-0.18, -0.8)$){$u^*_R$};
						\node[font=\footnotesize] at ($(\xU+\dR, \h) + ( 0.18, -0.8)$){$u^*_L$};
						
						%% --- campitura grigia dei due triangoli del campo u*
						\fill[black!12]
						(\xU, 0) -- (\xU-\dL, \h) -- (\xU, \h) -- cycle;
						\fill[black!12]
						(\xU, 0) -- (\xU+\dR, \h) -- (\xU, \h) -- cycle;
						
						%% chiusura orizzontale al top (alla quota h)
						\draw[line width=0.5pt](\xU-\dL, \h) -- (\xU+\dR, \h);
						
						\draw[line width=0.5pt](\xU, 0) -- (\xU, \h);
						
						\node[font=\small, align=center] at (1.5*\L, -3.40)
						{(b) Linea d'influenza di $N^{(4)}$};
						
					\end{scope}
					
					%% =========================================================================
					%% PANNELLO 3 — DISTORSIONE SULL'ASTA 2 (CORRENTE INFERIORE CENTRALE)
					%%              MECCANISMO VIRTUALE, CAMPI v* (sotto) E u* (a destra)
					%% =========================================================================
					\begin{scope}[xshift=0cm, yshift=\DY cm]
						
						%% --- coordinate dei nodi + centro C_R (fuori travatura)
						\coordinate (N1) at (\xNa, 0);
						\coordinate (N2) at (\xNb, 0);
						\coordinate (N3) at (\xNc, 0);
						\coordinate (N4) at (\xNd, 0);
						\coordinate (N5) at (\xNe, \h);
						\coordinate (N6) at (\xNf, \h);
						\coordinate (N7) at (\xNg, \h);
						\coordinate (CR) at (\xNd, 2*\h);
						
						%% --- riempimenti grigi dei due tronconi
						\fill[black!12] (N1) -- (N2) -- (N6) -- (N5) -- cycle;    % sinistro
						\fill[black!12] (N3) -- (N4) -- (N7) -- (N6) -- cycle;    % destro
						
						%% --- aste (asta 2, N2--N3, in grigio chiaro)
						\draw[asta](N1)--(N5);
						\draw[asta](N2)--(N5);
						\draw[asta](N2)--(N6);
						\draw[asta](N3)--(N6);
						\draw[asta](N3)--(N7);
						\draw[asta](N4)--(N7);
						\draw[asta](N1)--(N2);
						\draw[asta, gray!55](N2)--(N3);                            % asta 2: cede
						\draw[asta](N3)--(N4);
						\draw[asta](N5)--(N6)--(N7);
						
						%% --- distorsione sull'asta 2: due frecce antagoniste sotto l'asta
						\pgfmathsetmacro{\yf}{-0.30}
						\pgfmathsetmacro{\xfc}{0.5*(\xNb+\xNc)}
						\def\frL{0.30}
						\draw[->, >={Latex[length=1.6mm]}, line width=0.8pt]
						(\xfc-0.06, \yf) -- (\xfc-\frL, \yf);
						\draw[->, >={Latex[length=1.6mm]}, line width=0.8pt]
						(\xfc+0.06, \yf) -- (\xfc+\frL, \yf);
						\node[below, font=\footnotesize] at (\xfc, \yf-0.02){$\varrho_w=1$};
						
						%% --- cerniera in N1
						\draw[line width=1.0pt](N1)--++(-0.24,-0.40)--++(0.48,0)--cycle;
						\draw[line width=1.0pt]($(N1)+(-0.34,-0.40)$)--($(N1)+(0.34,-0.40)$);
						\foreach \x in {-0.28,-0.14,0.00,0.14,0.28}
						\draw[very thin]($(N1)+(\x,-0.40)$)--++(-.10,-0.13);
						
						%% --- carrello in N4
						\draw[line width=1.0pt](N4)--++(-0.24,-0.34)--++(0.48,0)--cycle;
						\draw[line width=1.0pt]($(N4)+(-0.34,-0.34)$)--($(N4)+(0.34,-0.34)$);
						\draw[line width=1.0pt]($(N4)+(-0.16,-0.47)$) circle(2.2pt);
						\draw[line width=1.0pt]($(N4)+( 0.16,-0.47)$) circle(2.2pt);
						\draw[line width=1.0pt]($(N4)+(-0.34,-0.58)$)--($(N4)+(0.34,-0.58)$);
						\foreach \x in {-0.28,-0.14,0.00,0.14,0.28}
						\draw[very thin]($(N4)+(\x,-0.58)$)--++(-.10,-0.13);
						
						%% --- linee di richiamo per C_R
						%%     verticale da N4 (carrello) + retta dei tre centri da N1 passante per N6
						\draw[ext, dashed, gray!70](N4) -- (CR);
						\draw[ext, dashed, gray!70](N1) -- (CR);
						
						%% --- pallini dei nodi e di C_R
						\foreach \n in {N1,N2,N3,N4,N5,N6,N7}
						{\node[nodo]at(\n){};}
						\node[nodo] at (CR){};
						
						%% --- etichette dei centri sulla struttura
						\node[above left,  font=\footnotesize] at (N1){$C_L$};
						\node[above,       font=\footnotesize] at (N6){$C_{L,R}$};
						\node[above right, font=\footnotesize] at (CR){$C_R$};
						
						%% ═════════════════════════════════════════════════════════════════
						%% RETTA DI PROIEZIONE ORIZZONTALE — campo v* (sotto la travatura)
						%% ═════════════════════════════════════════════════════════════════
						\def\yB{-2.30}
						\def\vpk{0.90}
						
						%% linee di richiamo verticali dai centri alla baseline
						\draw[ext, dashed, gray!70](\xNa, -0.78)--(\xNa, \yB);
						\draw[ext, dashed, gray!70](\xNf,  \h )--(\xNf, \yB);
						\draw[ext, dashed, gray!70](\xNd, -0.78)--(\xNd, \yB);
						
						%% area della LI (picco SOTTO la baseline)
						\fill[black!12]
						(\xNa, \yB) -- (\xNf, \yB-\vpk) -- (\xNd, \yB) -- cycle;
						%% spezzata
						\draw[line width=1.0pt]
						(\xNa, \yB) -- (\xNf, \yB-\vpk) -- (\xNd, \yB);
						%% baseline
						\draw[line width=0.4pt](\xNa-0.40, \yB)--(\xNd+0.40, \yB);
						
						%% tick sulle proiezioni dei centri
						\foreach \x in {\xNa, \xNf, \xNd}
						\draw[line width=0.5pt](\x, \yB-0.08)--(\x, \yB+0.08);
						
						%% labels dei centri (SOPRA la baseline: la LI sta sotto)
						\node[above, font=\footnotesize] at (\xNa, \yB+0.08){$C_L$};
						\node[above, font=\footnotesize] at (\xNf, \yB+0.08){$C_{L,R}$};
						\node[above, font=\footnotesize] at (\xNd, \yB+0.08){$C_R$};
						
						%% valore del picco (etichetta sotto la punta)
						\node[below, font=\footnotesize]
						at (\xNf, \yB-\vpk-0.03){$\dfrac{3}{4}\,\dfrac{\varrho_w\,\ell}{h}$};
						
						%% label campo v*
						%\node[left, font=\small] at (\xNa-0.45, \yB-\vpk/2){$v^*$};
						%% etichette dei due campi v*_L e v*_R (sotto la baseline, fuori dal triangolo)
						\node[below, font=\footnotesize]
						at ($(\xNa, \yB)!0.5!(\xNf, \yB-\vpk) + (0, -0.22)$){$v^*_L$};
						\node[below, font=\footnotesize]
						at ($(\xNf, \yB-\vpk)!0.5!(\xNd, \yB) + (0, -0.22)$){$v^*_R$};
						
						
						%% ═════════════════════════════════════════════════════════════════
						%% RETTA DI PROIEZIONE VERTICALE — campo u* (a destra della travatura)
						%% ═════════════════════════════════════════════════════════════════
						\pgfmathsetmacro{\xU}{\xNd+1.80}
						\def\uRho{0.80}
						\pgfmathsetmacro{\uHalf}{0.5*\uRho}
						
						%% linee di richiamo orizzontali dalla struttura alla retta u*
						\draw[ext, dashed, gray!70](\xNd+0.40,   0) -- (\xU+0.10,   0);
						\draw[ext, dashed, gray!70](\xNg+0.10,  \h) -- (\xU+0.10,  \h);
						\draw[ext, dashed, gray!70]      (CR)       -- (\xU+0.10, 2*\h);
						
						
						%% --- riempimento grigio: trapezio tra le rette u* (limitato a y ∈ [0, h])
						\fill[black!12]
						(\xU, 0) -- (\xU+\uRho, 0) -- (\xU+\uHalf, \h) -- (\xU, \h) -- cycle;
						
						%% chiusure orizzontali del trapezio (basi superiore e inferiore)
						\draw[line width=0.5pt](\xU, 0) -- (\xU+\uRho, 0);
						\draw[line width=0.5pt](\xU, \h) -- (\xU+\uHalf, \h);
						
						%			%% retta verticale di riferimento (tratteggiata) — ridisegnata sopra il fill
						%			\draw[ext, dashed, gray!70](\xU, 0) -- (\xU, 2*\h);
						
						%% retta verticale di riferimento: continua su y ∈ [0, h], tratteggiata su [h, 2h]
						\draw[line width=0.5pt](\xU, 0) -- (\xU, \h);
						\draw[ext, dashed, gray!70](\xU, \h) -- (\xU, 2*\h);
						
						%% tick e labels dei tre centri sulla retta di riferimento
						\draw[line width=0.5pt](\xU-0.08,   0 )--(\xU+0.08,   0 );
						\node[left, font=\footnotesize] at (\xU-0.10,   0 ){$C_L$};
						\draw[line width=0.5pt](\xU-0.08,  \h )--(\xU+0.08,  \h );
						\node[left, font=\footnotesize] at (\xU-0.10,  \h ){$C_{L,R}$};
						\draw[line width=0.5pt](\xU-0.08, 2*\h)--(\xU+0.08, 2*\h);
						\node[left, font=\footnotesize] at (\xU-0.10, 2*\h){$C_R$};
						
						%% u*_L e u*_R (tracciate sopra il fill)
						\draw[line width=0.9pt](\xU, 0) -- (\xU+\uHalf, \h);
						\draw[line width=0.9pt](\xU+\uRho, 0) -- (\xU, 2*\h);
						
						%% quota del gap ρ_w alla base
						\draw[ext](\xU,       0) -- (\xU,       -0.30);
						\draw[ext](\xU+\uRho, 0) -- (\xU+\uRho, -0.30);
						\draw[quota]
						(\xU,       -0.25) -- (\xU+\uRho, -0.25)
						node[midway, below, font=\footnotesize]{$\varrho_w$};
						
						%% labels dei due campi
						\node[right=1pt, font=\footnotesize] at (\xU+\uHalf-0.7, 0.75*\h){$u^*_L$};
						\node[right=1pt, font=\footnotesize] at (\xU+0.65*\uRho, 0.5*\h){$u^*_R$};
						
						\node[font=\small, align=center] at (1.5*\L, -4.80)
						{(c) Linea d'influenza di $N^{(2)}$};
					\end{scope}
					
					%% =========================================================================
					%% PANNELLO 4 — DISTORSIONE SULL'ASTA 8 (DIAGONALE N2-N6)
					%%              STRUTTURA CAMPITA (senza meccanismo — da completare)
					%% =========================================================================
					\begin{scope}[xshift=\DX cm, yshift=\DY cm]
						
						%% --- coordinate dei nodi
						\coordinate (N1) at (\xNa, 0);
						\coordinate (N2) at (\xNb, 0);
						\coordinate (N3) at (\xNc, 0);
						\coordinate (N4) at (\xNd, 0);
						\coordinate (N5) at (\xNe, \h);
						\coordinate (N6) at (\xNf, \h);
						\coordinate (N7) at (\xNg, \h);
						
						%% --- riempimenti grigi dei tronconi
						%% sinistro: triangolo N1-N2-N5
						\fill[black!12] (N1) -- (N2) -- (N5) -- cycle;
						%% destro: poligono N3-N4-N7-N6 (due triangoli adiacenti uniti)
						\fill[black!12] (N3) -- (N4) -- (N7) -- (N6) -- cycle;
						
						%% --- aste
						%% aste del troncone sinistro (nere)
						\draw[asta](N1)--(N2);            % asta 1
						\draw[asta](N1)--(N5);            % asta 6
						\draw[asta](N2)--(N5);            % asta 7
						%% aste pendenti (nere, tratto normale) — collegano i due tronconi
						\draw[asta](N2)--(N3);            % asta 2
						\draw[asta](N5)--(N6);            % asta 4
						%% asta distorta (grigio chiaro)
						\draw[asta, gray!55](N2)--(N6);   % asta 8: cede
						%% aste del troncone destro (nere)
						\draw[asta](N3)--(N6);            % asta 9
						\draw[asta](N3)--(N7);            % asta 10
						\draw[asta](N3)--(N4);            % asta 3
						\draw[asta](N6)--(N7);            % asta 5
						\draw[asta](N4)--(N7);            % asta 11
						
						%% --- distorsione sull'asta 8: due frecce antagoniste accanto all'asta
						%%     (l'asta è diagonale, le frecce le allineiamo parallele all'asta)
						\pgfmathsetmacro{\xmid}{0.5*(\xNb+\xNf)}    % x medio asta 8
						\pgfmathsetmacro{\ymid}{0.5*\h}             % y medio asta 8
						%% versore lungo l'asta (N2 -> N6): dx = xNf-xNb = 0.5*ell, dy = h
						\pgfmathsetmacro{\dAx}{\xNf-\xNb}
						\pgfmathsetmacro{\dAy}{\h}
						\pgfmathsetmacro{\Lmod}{sqrt(\dAx*\dAx + \dAy*\dAy)}
						\pgfmathsetmacro{\tux}{\dAx/\Lmod}           % tangente unitaria x
						\pgfmathsetmacro{\tuy}{\dAy/\Lmod}           % tangente unitaria y
						\pgfmathsetmacro{\nux}{-\tuy}                % normale unitaria x (a sinistra)
						\pgfmathsetmacro{\nuy}{ \tux}                % normale unitaria y
						\def\off{0.22}                               % offset laterale delle frecce
						\def\frL{0.35}                               % lunghezza di ciascuna freccia
						\def\gap{0.06}                               % gap centrale tra le frecce
						%% punto base delle due frecce (traslato lateralmente dall'asta)
						\pgfmathsetmacro{\Bx}{\xmid + \off*\nux}
						\pgfmathsetmacro{\By}{\ymid + \off*\nuy}
						%% freccia verso N2 (giù-sinistra)
						\draw[->, >={Latex[length=1.6mm]}, line width=0.8pt]
						(\Bx - \gap*\tux, \By - \gap*\tuy) --
						(\Bx - \frL*\tux, \By - \frL*\tuy);
						%% freccia verso N6 (su-destra)
						\draw[->, >={Latex[length=1.6mm]}, line width=0.8pt]
						(\Bx + \gap*\tux, \By + \gap*\tuy) --
						(\Bx + \frL*\tux, \By + \frL*\tuy);
						%% etichetta ρ_w al centro (spostata dalla parte opposta alle frecce)
						\node[font=\footnotesize]
						at (\Bx + 0.20*\nux, \By + 0.20*\nuy){$\varrho_w$};
						
						%% --- cerniera in N1
						\draw[line width=1.0pt](N1)--++(-0.24,-0.40)--++(0.48,0)--cycle;
						\draw[line width=1.0pt]($(N1)+(-0.34,-0.40)$)--($(N1)+(0.34,-0.40)$);
						\foreach \x in {-0.28,-0.14,0.00,0.14,0.28}
						\draw[very thin]($(N1)+(\x,-0.40)$)--++(-.10,-0.13);
						
						%% --- carrello in N4
						\draw[line width=1.0pt](N4)--++(-0.24,-0.34)--++(0.48,0)--cycle;
						\draw[line width=1.0pt]($(N4)+(-0.34,-0.34)$)--($(N4)+(0.34,-0.34)$);
						\draw[line width=1.0pt]($(N4)+(-0.16,-0.47)$) circle(2.2pt);
						\draw[line width=1.0pt]($(N4)+( 0.16,-0.47)$) circle(2.2pt);
						\draw[line width=1.0pt]($(N4)+(-0.34,-0.58)$)--($(N4)+(0.34,-0.58)$);
						\foreach \x in {-0.28,-0.14,0.00,0.14,0.28}
						\draw[very thin]($(N4)+(\x,-0.58)$)--++(-.10,-0.13);
						
						%% --- pallini dei nodi
						\foreach \n in {N1,N2,N3,N4,N5,N6,N7}
						{\node[nodo]at(\n){};}
						
						%% ═════════════════════════════════════════════════════════════════
						%% RETTA DI PROIEZIONE ORIZZONTALE — campo v* (attraversa la baseline)
						%% Ramo sinistro SOTTO la baseline (N2 scende)
						%% Ramo destro SOPRA la baseline (N3 sale)
						%% Raccordo rettilineo tra (ell, -vpk) e (2ell, +vpk)
						%% ═════════════════════════════════════════════════════════════════
						\def\yB{-2.30}
						\def\vpk{0.60}
						\pgfmathsetmacro{\xcross}{0.5*(\xNb+\xNc)}    % ascissa dove il raccordo taglia la baseline
						
						%% linee di richiamo verticali dai punti chiave
						\draw[ext, dashed, gray!70](\xNa, -0.78)--(\xNa, \yB);
						\draw[ext, dashed, gray!70](\xNb,  0   )--(\xNb, \yB-\vpk);
						\draw[ext, dashed, gray!70](\xNf, \h   )--(\xNf, \yB+1.5*\vpk);
						\draw[ext, dashed, gray!70](\xNd, -0.78)--(\xNd, \yB);
						
						
						%% campitura: lobo sinistro (sotto baseline), esteso da C_L a proiezione di N2
						\fill[black!12]
						(\xNa, \yB) -- (\xNb, \yB-\vpk) -- (\xNb, \yB) -- cycle;
						%% campitura: lobo destro (sopra baseline), esteso da proiezione di N6 a C_R
						\fill[black!12]
						(\xNf, \yB) -- (\xNf, \yB+1.5*\vpk) -- (\xNd, \yB) -- cycle;
						
						%% ramo sinistro: da C_L a N2
						\draw[line width=1.0pt]
						(\xNa, \yB) -- (\xNb, \yB-\vpk);
						%% ramo destro: da proiezione N6 a C_R
						\draw[line width=1.0pt]
						(\xNf, \yB+1.5*\vpk) -- (\xNd, \yB);
						
						%					%% spezzata della LI (ramo sinistro + raccordo + ramo destro)
						%					\draw[line width=1.0pt]
						%					(\xNa, \yB) -- (\xNb, \yB-\vpk)
						%					-- (\xNc, \yB+\vpk) -- (\xNd, \yB);
						
						%% baseline
						\draw[line width=0.4pt](\xNa-0.40, \yB)--(\xNd+0.40, \yB);
						
						%% tick solo sui due centri propri
						\foreach \x in {\xNa, \xNd}
						\draw[line width=0.5pt](\x, \yB-0.08)--(\x, \yB+0.08);
						
						%% labels dei centri sulla baseline
						\node[below, font=\footnotesize] at (\xNa, \yB-0.08){$C_L$};
						\node[below, font=\footnotesize] at (\xNd, \yB-0.08){$C_R$};
						
						
						%% linee sottili di chiusura verticale dei due campi
						\draw[line width=0.5pt] (\xNb, \yB) -- (\xNb, \yB-\vpk);       % chiude v*_L a destra
						\draw[line width=0.5pt] (\xNf, \yB) -- (\xNf, \yB+1.5*\vpk);   % chiude v*_R a sinistra
						
						
						%% valori dei picchi
						%					\node[below, font=\footnotesize]
						%					at (\xNb, \yB-\vpk-0.03){$\vartheta\,\ell$};
						%					\node[above, font=\footnotesize]
						%					at (\xNc, \yB+\vpk+0.03){$\vartheta\,\ell$};
						\node[below, font=\footnotesize]
						at (\xNb, \yB-\vpk-0.03){$\vartheta^{*}\ell$};
						\node[above, font=\footnotesize]
						at (\xNf, \yB+1.5*\vpk+0.03){$\tfrac{3}{2}\,\vartheta^{*}\ell$};
						
						
						
						%% etichette dei due rami
						\node[below, font=\footnotesize]
						at ($(\xNa, \yB)!0.5!(\xNb, \yB-\vpk) + (0, -0.22)$){$v^*_L$};
						\node[above, font=\footnotesize]
						at ($(\xNc, \yB+\vpk)!0.5!(\xNd-0.25, \yB) + (0, 0.22)$){$v^*_R$};
						
						
						
						%% ═════════════════════════════════════════════════════════════════
						%% RETTA DI PROIEZIONE VERTICALE — campo u* (a destra della travatura)
						%% u*_L ≡ u*_R : singola retta, singolo triangolo di campitura
						%% ═════════════════════════════════════════════════════════════════
						\pgfmathsetmacro{\xU}{\xNd+1.40}
						\def\uRho{0.60}        % valore visuale di vartheta*h al top
						
						%% linee di richiamo orizzontali dalla struttura alla retta u*
						\draw[ext, dashed, gray!70](\xNd+0.40, 0 )--(\xU+0.10, 0 );
						\draw[ext, dashed, gray!70](\xNg+0.10, \h)--(\xU+0.10, \h);
						
						%% campitura del triangolo unico
						\fill[black!12]
						(\xU, 0) -- (\xU+\uRho, \h) -- (\xU, \h) -- cycle;
						
						%% retta verticale di riferimento: continua su [0, h]
						\draw[line width=0.5pt](\xU, 0) -- (\xU, \h);
						
						%% chiusure orizzontali del triangolo
						\draw[line width=0.5pt](\xU, 0) -- (\xU, 0);           % base: è il punto C_L≡C_R
						\draw[line width=0.5pt](\xU, \h) -- (\xU+\uRho, \h);   % top
						
						%% tick e labels dei due centri sovrapposti al piede e di C_R... wait
						%% in realtà C_L e C_R stanno ENTRAMBI a y=0, sovrapposti
						\draw[line width=0.5pt](\xU-0.08, 0 )--(\xU+0.08, 0 );
						\node[left, font=\footnotesize] at (\xU-0.10, 0){$C_L\!\equiv\!C_R$};
						
						%% retta u*_L ≡ u*_R
						\draw[line width=0.9pt](\xU, 0) -- (\xU+\uRho, \h);
						
						%					%% archetto della rotazione ϑ* al piede (unico, perché ϑ_L = ϑ_R)
						%					\def\Rau{0.55}
						%					\pgfmathsetmacro{\angU}{atan2(\h, \uRho)}
						%					\draw[line width=0.5pt]
						%					(\xU, 0) ++(\angU:\Rau) arc (\angU:90:\Rau);
						%					\node[font=\footnotesize]
						%					at ($(\xU, 0) + ({0.5*(\angU+90)}:\Rau) + (0.22, 0.05)$){$\vartheta^*$};
						
						%% valore dell'ordinata al top
						\node[right=1pt, font=\footnotesize] at (\xU+\uRho-0.8, \h+0.3){$\vartheta^*\,h$};
						
						%% label del campo
						\node[font=\footnotesize]
						at ($(\xU, 0)!0.5!(\xU+\uRho, \h) + (0.72, -0.05)$){$u^*_L\!\equiv\!u^*_R$};
						
						%% --- piedi delle perpendicolari da C_L e C_R sulla retta asta 8
						\coordinate (HL) at ($(N2)!(N1)!(N6)$);      % piede da C_L = N1
						\coordinate (HR) at ($(N2)!(N4)!(N6)$);      % piede da C_R = N4
						
						%% --- prolungamento (tratteggiato) della retta asta 8 fino ai due piedi
						\draw[ext, dashed, gray!70] (HL) -- (N2);
						\draw[ext, dashed, gray!70] (N6) -- (HR);
						
						%% --- segmenti d_L e d_R a tratto continuo
						\draw[line width=0.4pt] (N1) -- (HL);
						\draw[line width=0.4pt] (N4) -- (HR);
						
						%% --- squadrette di ortogonalità ai piedi
						\def\sq{0.12}
						\pgfmathsetmacro{\Lmod}{sqrt((\xNf-\xNb)^2 + \h*\h)}
						\pgfmathsetmacro{\tux}{(\xNf-\xNb)/\Lmod}
						\pgfmathsetmacro{\tuy}{\h/\Lmod}
						\pgfmathsetmacro{\nux}{-\tuy}
						\pgfmathsetmacro{\nuy}{ \tux}
						
						%% squadretta in HL: tra segmento (HL -> N1) e prolungamento (HL -> N2)
						%% N1 sta dal lato -n rispetto alla retta; N2 sta in direzione +t da HL.
						\draw[line width=0.3pt]
						($(HL) + (-\sq*\nux, -\sq*\nuy)$) --
						($(HL) + (-\sq*\nux + \sq*\tux, -\sq*\nuy + \sq*\tuy)$) --
						($(HL) + ( \sq*\tux,  \sq*\tuy)$);
						
						%% squadretta in HR: tra segmento (HR -> N4) e prolungamento (HR -> N6)
						%% N4 sta dal lato -n; N6 sta in direzione -t da HR.
						\draw[line width=0.3pt]
						($(HR) + (-\sq*\nux, -\sq*\nuy)$) --
						($(HR) + (-\sq*\nux - \sq*\tux, -\sq*\nuy - \sq*\tuy)$) --
						($(HR) + (-\sq*\tux, -\sq*\tuy)$);
						
						%% --- labels d_L e d_R a metà dei segmenti
						\node[font=\footnotesize]
						at ($(N1)!0.5!(HL) + (0, -0.18)$){$d_L$};
						\node[font=\footnotesize]
						at ($(N4)!0.5!(HR) + (-0.40, 0.05)$){$d_R$};
						
						\node[font=\small, align=center] at (1.5*\L, -4.80)
						{(d) Linea d'influenza di $N^{(8)}$};
					\end{scope}
					
				\end{tikzpicture}%
			}
			\caption{Travatura reticolare piana a correnti paralleli: linee d'influenza
				degli sforzi normali nelle aste~$4$ (corrente superiore),
				$2$ (corrente inferiore) e~$8$ (diagonale), sotto carico unitario
				verticale mobile sui nodi dei due correnti.
				\emph{(a)} Geometria della travatura, con luce totale $3\ell$ e altezza
				$h$ fra i correnti.
				\emph{(b)}, \emph{(c)}, \emph{(d)} Meccanismi virtuali duali (TSV)
				ottenuti imponendo una distorsione longitudinale $\varrho_w$ sull'asta
				in esame; il campo $v^*$ sotto la travatura fornisce le ordinate della
				linea d'influenza per carichi verticali, il campo $u^*$ a fianco le
				ordinate per carichi orizzontali.}
			\label{fig:LI_reticolare}
		\end{figure}
